\documentclass[12pt]{article}

\usepackage[brazilian]{babel}
\usepackage[utf8]{inputenc}
\usepackage[T1]{fontenc}
\usepackage{amsmath,amssymb}
\usepackage{graphicx}
\usepackage{geometry}
\usepackage{xcolor}  % ou \usepackage{color}

\geometry{margin=2.5cm}

\title{Empacotamento Tridimensional de Cargas com Restrições de Estabilidade e Ordem de Entrega Utilizando BRKGA}
\author{Danillo Mendes Santiago \ 10414592}
\date{Junho de 2026}

\begin{document}

\maketitle

\begin{abstract}
Este trabalho aborda o problema de empacotamento tridimensional de caixas em um contêiner. Além das restrições geométricas usuais, são consideradas restrições de estabilidade física e de ordem de descarregamento. O problema é modelado como um problema de otimização combinatória e resolvido por meio de um algoritmo Biased Random-Key Genetic Algorithm (BRKGA).
\end{abstract}

\section{Introdução}

Problemas de empacotamento aparecem frequentemente em logística, transporte e armazenagem. O objetivo consiste em posicionar objetos dentro de um espaço limitado utilizando o volume disponível de forma eficiente.

Neste trabalho considera-se um conjunto de caixas retangulares com dimensões distintas que devem ser acomodadas em um contêiner retangular. Além do aproveitamento de espaço, devem ser respeitadas restrições de estabilidade e de ordem de entrega.

\section{Descrição do Problema}

Seja um contêiner de dimensões

\[
(L,W,H).
\]

Considere um conjunto de caixas

\[
B=\{B_1,B_2,\ldots,B_n\},
\]

onde cada caixa $B_i$ possui dimensões

\[
(l_i,w_i,h_i).
\]

Cada caixa possui também uma prioridade de entrega

\[
d_i \in \mathbb{N}.
\]

Valores menores de $d_i$ indicam entregas realizadas antes.

A face localizada em

\[
x=0
\]

é definida como porta de descarregamento.

O objetivo é determinar uma disposição das caixas que maximize a ocupação do contêiner e respeite todas as restrições do problema.

\section{Modelagem}

O espaço do contêiner é discretizado em uma grade de vóxels, onde cada coordenada $(x,y,z) \in \mathbb{Z}^3$ representa uma célula unitária.

Cada caixa $B_i$ é representada pela tupla

\[
\text{caixa}_i = (x_i, y_i, z_i, r_i, d_i) \in \mathbb{Z}^3 \times \{0,1\} \times \{0,1,2,\ldots,D-1\},
\]

onde:
\begin{itemize}
\item $(x_i, y_i, z_i)$ são as coordenadas discretizadas do posicionamento da caixa;
\item $r_i \in \{0,1\}$ indica a orientação: $r_i = 0$ mantém as dimensões originais $(l_i, w_i, h_i)$; $r_i = 1$ aplica uma rotação de $90^\circ$ no plano horizontal, trocando $l_i$ com $w_i$;
\item $d_i$ é a prioridade de entrega.
\end{itemize}


\colorbox{yellow}{%
    \begin{minipage}{0.9\linewidth}
        Justificar cada um:
        \begin{itemize}
            \item discretização e posição, num subconjunto de Z³.
            \item Devido à simetria do paralelepípedo retângulo, duas rotações de $90^\circ$ recuperam as dimensões originais. Portanto, apenas duas orientações são consideradas.
            \item talvez já esteja claro, explicar que $D=\#\textit{pedidos}$.
        \end{itemize}
    \end{minipage}%
}



A solução deve satisfazer as seguintes condições.

\subsection{Permanência no Contêiner}

Toda caixa deve permanecer integralmente dentro do contêiner. Seja $(\tilde{l}_i, \tilde{w}_i, \tilde{h}_i)$ as dimensões efetivas após aplicação da rotação $r_i$:

\[
(\tilde{l}_i, \tilde{w}_i, \tilde{h}_i) = 
\begin{cases}
(l_i, w_i, h_i), & \text{se } r_i = 0, \\
(w_i, l_i, h_i), & \text{se } r_i = 1.
\end{cases}
\]

As coordenadas devem satisfazer:

\[
0 \leq x_i \leq L - \tilde{l}_i,
\]
\[
0 \leq y_i \leq W - \tilde{w}_i,
\]
\[
0 \leq z_i \leq H - \tilde{h}_i.
\]

\subsection{Não Sobreposição}

Duas caixas não podem ocupar simultaneamente a mesma região do espaço. Para quaisquer caixas distintas $B_i$ e $B_j$, os intervalos ocupados em cada eixo não devem se sobrepor simultaneamente:

\[
[x_i, x_i+\tilde{l}_i) \cap [x_j, x_j+\tilde{l}_j) = \emptyset \quad \text{ou} \quad [y_i, y_i+\tilde{w}_i) \cap [y_j, y_j+\tilde{w}_j) = \emptyset \quad \text{ou} \quad [z_i, z_i+\tilde{h}_i) \cap [z_j, z_j+\tilde{h}_j) = \emptyset.
\]

\subsection{Estabilidade}

Nenhuma caixa pode permanecer flutuando. Uma caixa com $z_i > 0$ é considerada estável se existir suporte contínuo em sua projeção horizontal:

\[
\forall (x,y) \in [x_i, x_i+\tilde{l}_i) \times [y_i, y_i+\tilde{w}_i), \quad \exists B_j \neq B_i \text{ tal que } z_j + \tilde{h}_j = z_i.
\]

Além disso, a projeção vertical do centro de massa da caixa deve pertencer à região de apoio, reduzindo o risco de tombamento durante a movimentação do contêiner.

\colorbox{yellow}{
    \begin{minipage}{0.9\linewidth}
        Aqui deve explicar melhor a projeção estar contida totalmente na caixa de baixo
    \end{minipage}
}


\subsection{Ordem de Entrega}

As caixas destinadas às primeiras entregas devem permanecer mais próximas da porta. Para quaisquer duas caixas $B_i$ e $B_j$ com $d_i < d_j$, deseja-se que

\[
x_i \leq x_j.
\]

Caixas de entregas posteriores não devem bloquear o acesso às caixas de entregas anteriores.

\section{Função Objetivo}

A qualidade de uma solução é avaliada por três critérios:

\begin{enumerate}
\item aproveitamento de volume;
\item estabilidade;
\item facilidade de descarregamento.
\end{enumerate}

A função de aptidão é definida por

\[
f = \alpha U - \beta P_s - \gamma P_d,
\]

onde

\begin{itemize}
\item $U$ representa o volume efetivamente utilizado;
\item $P_s$ representa penalizações por instabilidade;
\item $P_d$ representa penalizações por violações da ordem de entrega;
\item $\alpha,\beta,\gamma > 0$ representam a importância de cada critério.
\end{itemize}

Soluções inviáveis recebem penalizações elevadas.

\colorbox{yellow}{
    \begin{minipage}{0.9\linewidth}
        Aqui deve detalhar mais um pouco, falar de maximizar f, etc
    \end{minipage}
}

\section{Representação no BRKGA}

Cada indivíduo é representado por um vetor de chaves aleatórias

\[
\mathbf{r} = (r_1, r_2, \ldots, r_m),
\]

com

\[
r_k \in [0,1].
\]

O decodificador interpreta essas chaves para definir:

\begin{enumerate}
\item ordem de inserção das caixas;
\item orientação de cada caixa;
\item posição de empacotamento.
\end{enumerate}

A partir dessas informações é construída uma solução completa.

\colorbox{yellow}{
    \begin{minipage}{0.9\linewidth}
        Como r representa uma solução? Como garantir que sempre será um solução válida? Como é montado r com 3 caixas? Exemplos, etc
    \end{minipage}
}

\section{Decodificador}

O decodificador executa as seguintes etapas:

\begin{enumerate}
\item ordenar as caixas segundo as chaves aleatórias;
\item determinar a orientação $r_i$ de cada caixa;
\item inserir cada caixa em uma posição viável $(x_i, y_i, z_i)$;
\item verificar sobreposição;
\item verificar estabilidade;
\item calcular penalizações;
\item avaliar a função de aptidão.
\end{enumerate}

O decodificador é responsável por transformar um cromossomo em uma solução do problema.

\section{Preparo, NFP, IFP, etc}


\section{Algoritmo BRKGA}

O método utilizado segue o fluxo:

\begin{enumerate}
\item gerar população inicial;
\item avaliar indivíduos;
\item selecionar elite;
\item gerar descendentes por cruzamento enviesado;
\item inserir mutantes aleatórios;
\item formar nova população;
\item repetir até atingir o critério de parada.
\end{enumerate}

A cada geração é armazenada a melhor solução encontrada.

\section{Implementação}

A implementação foi desenvolvida em Python.

As principais estruturas utilizadas são:

\begin{itemize}
\item classe para representar caixas, armazenando dimensões $(l_i, w_i, h_i)$ e prioridade $d_i$;
\item classe para representar o contêiner com dimensões $(L, W, H)$;
\item classe para indivíduos do BRKGA;
\item decodificador responsável pelo empacotamento;
\item função de avaliação.
\end{itemize}

A execução do algoritmo produz uma disposição das caixas dentro do contêiner e o valor da função objetivo associado.

\section{Resultados}

Nesta seção serão apresentados:

\begin{itemize}
\item instâncias utilizadas;
\item parâmetros do BRKGA;
\item taxa de ocupação obtida;
\item análise da estabilidade;
\item análise da ordem de descarregamento.
\end{itemize}

\section{Conclusão}

Foi apresentada uma modelagem para o problema de empacotamento tridimensional com restrições de estabilidade e logística de entrega. A utilização do BRKGA permite explorar eficientemente o espaço de soluções, produzindo arranjos viáveis para problemas de grande porte.

\bibliographystyle{plain}
\bibliography{referencias}

\end{document}